\chapter{The lower-bound reduction}

Fix a horizon $T\geq 1$ and a predetermined nonnegative step-size schedule
$h=(h_t)_{t<T}$.  The formal development associates to $h$ a finite scalar
functional $\mathcal F(h)$.  The proof has two largely independent parts:
first realize $\mathcal F(h)$ by a smooth convex gradient-descent instance,
then prove a schedule-independent lower bound for $\mathcal F(h)$.

The target exponent is
\[
  p_0=\frac{45720}{25000}=1.8288,
  \qquad
  \beta_0=p_0-1=\frac{20720}{25000}=0.8288.
\]

\begin{definition}[Normalized schedule floor]
  \label{def:normalized-floor}
  \lean{GDLowerBound.RankAnalysis.NormalizedScheduleFloor}
  \leanok
  A normalized schedule floor at exponent $p$ and constant $c$ is the
  assertion that every nonnegative schedule satisfies
  \[
    \frac{c}{B(h)(r(h)+1)^{p-1}}\leq \mathcal F(h),
  \]
  where $B(h)$ is the capped schedule mass and $r(h)$ is the number of long
  steps.
\end{definition}

\begin{theorem}[Normalized-floor reduction]
  \label{thm:floor-reduction}
  \lean{GDLowerBound.RankAnalysis.functionalFloor_of_normalizedScheduleFloor}
  \leanok
  A normalized schedule floor at exponent $p\geq1$ implies the corresponding
  horizon-rate functional floor, with only the explicit factor two used by
  the geometric realization.
\end{theorem}

\begin{proof}
  \leanok
  \uses{def:normalized-floor}
  Compare capped mass and rank with the horizon, then rewrite the factor-two
  normalization of the objective gap.
\end{proof}

\section{Geometric realization}

\begin{theorem}[Chronological-chain realization]
  \label{thm:geometric-realization}
  \lean{GDLowerBound.GeometricRealization.geometricRealization}
  \leanok
  Every admissibly weighted chronological chain is realized by a
  $1$-smooth convex objective and an exact gradient-descent trajectory.  The
  initial point has distance one from the minimizer, and the terminal
  objective gap equals the chain product divided by twice its terminal
  scale.
\end{theorem}

\begin{proof}
  \leanok
  The block vectors, convex projection envelope, and saturated amplitudes
  give the objective and trajectory explicitly.  The empty-chain case is
  included in the same construction.
\end{proof}

\begin{corollary}[Functional attainment]
  \label{cor:functional-attainment}
  \lean{GDLowerBound.functionalAttainment}
  \leanok
  For every schedule, a maximizing chronological chain yields a normalized
  smooth convex instance in a dimension between $1$ and $T+1$ whose terminal
  objective gap is exactly $\mathcal F(h)/2$.
\end{corollary}

\begin{proof}
  \leanok
  \uses{thm:geometric-realization}
  Choose a maximizing chain and apply its geometric realization.
\end{proof}

\section{Rank propagation and the exact envelope}

\begin{lemma}[One-step Lyapunov inequality]
  \label{lem:one-step-lyapunov}
  \lean{GDLowerBound.RankAnalysis.oneStepLyapunov}
  \leanok
  Adjacent reciprocal-mass states satisfy a Lyapunov increment estimate with
  an explicit $O(n^{-2})$ discretization remainder.
\end{lemma}

\begin{proof}
  \leanok
  The adjacent-rank mass and state recurrences are substituted into the
  rational Lyapunov potential and the resulting scalar inequality is
  certified exactly.
\end{proof}

\begin{theorem}[Boundary propagation]
  \label{thm:boundary-propagation}
  \lean{GDLowerBound.RankAnalysis.boundaryPropagation}
  \leanok
  On a rank interval satisfying the cutoff conditions, unresolved mass at
  the initial rank is controlled by terminal capped mass and the prescribed
  power of the rank ratio.
\end{theorem}

\begin{proof}
  \leanok
  \uses{lem:one-step-lyapunov}
  Sum the one-step Lyapunov inequalities, telescope the potential, and retain
  the explicit boundary constant.
\end{proof}

\begin{theorem}[Bounded-rank prefix]
  \label{thm:bounded-rank}
  \lean{GDLowerBound.RankAnalysis.boundedRankPrefix_uniform}
  \leanok
  Every prefix of bounded rank contributes a fixed positive multiple of the
  reciprocal unresolved mass to $\mathcal F(h)$.
\end{theorem}

\begin{proof}
  \leanok
  A dense final rank is handled by chronological matching; a diffuse final
  rank is removed and the argument proceeds inductively.
\end{proof}

\begin{theorem}[Exact-envelope small-state branch]
  \label{thm:exact-envelope-small-state}
  \lean{GDLowerBound.FourBlock.chronologicalEdgeKernelProduct_lt_one_of_zeta_le,GDLowerBound.FourBlock.quarter_div_unresolvedMass_lt_functional_of_zeta_le}
  \leanok
  Above the fixed all-rank cutoff, a sufficiently small reciprocal-mass state
  makes the exact chronological edge-kernel product strictly less than one
  and forces
  \[
    \frac{1}{4D_q}<\mathcal F(h).
  \]
\end{theorem}

\begin{proof}
  \leanok
  Exact equalization and chronological outside-in matching bound the kernel
  product; the maximizing-chain formula then gives the functional branch.
\end{proof}

\begin{theorem}[Exact-envelope normalized floor]
  \label{thm:exact-envelope-floor}
  \lean{GDLowerBound.FourBlock.exactEnvelopeNormalizedLowerBound}
  \leanok
  At the rational exponent
  $1.8565576222695287$, there is a positive constant giving an unconditional
  normalized schedule floor.
\end{theorem}

\begin{proof}
  \leanok
  \uses{thm:boundary-propagation,thm:bounded-rank,thm:exact-envelope-small-state}
  Split at the last small state.  Use the functional branch there and
  propagate the remaining tail; if there is no small state, use the bounded
  prefix and propagate from the all-rank cutoff.
\end{proof}

\chapter{The four-block improvement}

The sharper route retains adjacent reciprocal-mass defects across blocks of
lengths $m$, $2m$, $3m$, and $4m$.  A certified local energy gap is combined
with finite harmonic averaging and a global maximal-dyadic scan.

\section{Rigidity and finite harmonic averaging}

\begin{lemma}[Sharp drift rigidity]
  \label{lem:sharp-rigidity}
  \lean{GDLowerBound.FourBlock.scheduleSharpRigidityBudget}
  \leanok
  The defect-weighted first moment of the relative mass increments is bounded
  by a telescoping critical potential plus an explicit reciprocal-square
  error.
\end{lemma}

\begin{proof}
  \leanok
  Sum the certified one-step sharp drift estimate over the rank interval.
\end{proof}

\begin{lemma}[Moment rigidity]
  \label{lem:moment-rigidity}
  \lean{GDLowerBound.FourBlock.scheduleMomentRigidityBudget}
  \leanok
  The squared deviation from the critical mean, together with the endpoint
  defect, is controlled by the mean deficit, a telescoping defect term, and
  an explicit reciprocal-square error.
\end{lemma}

\begin{proof}
  \leanok
  Sum the exact joint moment--defect one-step inequality.
\end{proof}

\begin{lemma}[Fixed-dilation sampling]
  \label{lem:fixed-dilation-sampling}
  \lean{GDLowerBound.FourBlock.fixedDilationHarmonicSample_le}
  \leanok
  Sampling a rank-normalized sequence at a fixed dilation $j$ is bounded by
  its full harmonic sum plus $j$ times its adjacent total variation.
\end{lemma}

\begin{proof}
  \leanok
  Partition the full rank interval into consecutive backward blocks and
  compare each selected endpoint with its block average.
\end{proof}

\begin{lemma}[Finite dilation-defect expansion]
  \label{lem:dilation-defect}
  \lean{GDLowerBound.FourBlock.dilationBlockH_upper}
  \leanok
  Every finite dilation defect is bounded by its middle-block mean deficit,
  its initial endpoint defect, a harmonic discretization term, and a
  reciprocal-square remainder.
\end{lemma}

\begin{proof}
  \leanok
  Expand the log mass ratio into adjacent increments and use the directed
  logarithm and quadratic certificates.
\end{proof}

\begin{lemma}[Middle-block occupancy identity]
  \label{lem:middle-block-occupancy}
  \lean{GDLowerBound.FourBlock.sum_middleBlocks_eq_occupancy}
  \leanok
  The harmonic average of the moving middle blocks equals a single rank sum
  weighted by the exact occupancy of each rank.
\end{lemma}

\begin{proof}
  \leanok
  Exchange two finite sums and identify the indicator sum with the occupancy
  weight.
\end{proof}

\begin{theorem}[Averaged schedule-energy bound]
  \label{thm:averaged-energy}
  \lean{GDLowerBound.FourBlock.averagedScheduleEnergy_upper}
  \leanok
  The harmonically averaged four-block schedule energy is bounded by the
  common square deviation, common defect, common mean deficit, and one named
  finite-error term.
\end{theorem}

\begin{proof}
  \leanok
  \uses{lem:sharp-rigidity,lem:moment-rigidity,lem:fixed-dilation-sampling,lem:dilation-defect,lem:middle-block-occupancy}
  Combine the endpoint, two-block, and three-block estimates; rewrite moving
  sums by occupancy; then apply the two common rigidity budgets.
\end{proof}

\begin{theorem}[Normalized local-gap comparison]
  \label{thm:normalized-local-gap}
  \lean{GDLowerBound.FourBlock.normalizedAveragedScheduleEnergy_lt_robustLocalGap}
  \leanok
  If the effective mean growth reaches $\beta_0$ and the complete finite
  error fits inside the certified margin, then the normalized averaged
  schedule energy is strictly below the robust local energy gap.
\end{theorem}

\begin{proof}
  \leanok
  \uses{thm:averaged-energy}
  Divide every term by the common positive harmonic weight, insert the
  rigidity budgets, and apply the exact rational local-gap certificate.
\end{proof}

\begin{corollary}[A good sampled scale]
  \label{cor:good-scale}
  \lean{GDLowerBound.FourBlock.normalizedAveraging_forces_small_scale}
  \leanok
  Under the same growth and error hypotheses, some sampled scale $m$ lies on
  the functional branch:
  \[
    \frac{1}{4D_{4m}}<\mathcal F(h).
  \]
\end{corollary}

\begin{proof}
  \leanok
  \uses{thm:normalized-local-gap}
  If every sampled scale were on the opposite side of the pointwise
  dichotomy, their positive weighted average would contradict the strict
  averaged local-gap bound.
\end{proof}

\section{Uniform dyadic windows and the global scan}

\begin{theorem}[Uniform finite window]
  \label{thm:uniform-window}
  \lean{GDLowerBound.FourBlock.exists_uniform_dyadicFourBlockWindow}
  \leanok
  There are fixed depth and base-scale thresholds beyond which every dyadic
  window has total normalized finite error strictly smaller than the local
  certificate margin.
\end{theorem}

\begin{proof}
  \leanok
  The reciprocal-square, endpoint-sampling, occupancy-boundary, and harmonic
  errors admit an explicit bound of the form $O(R^{-1})+O(R/M)$; choose the
  depth and base scale by the Archimedean property.
\end{proof}

\begin{theorem}[Dyadic window dichotomy]
  \label{thm:dyadic-dichotomy}
  \lean{GDLowerBound.FourBlock.finiteWindowDyadic_functional_or_massRatio_lt,GDLowerBound.FourBlock.exists_uniformFiniteWindowDyadicDichotomy}
  \leanok
  Every sufficiently large admissible dyadic window either contains a later
  functional rank or its endpoint unresolved-mass ratio grows with the exact
  coefficient $\beta_0$ in front of $R\log 2$, with explicit boundary terms.
\end{theorem}

\begin{proof}
  \leanok
  \uses{cor:good-scale,thm:uniform-window}
  If effective mean growth is at least $\beta_0$, extract a good scale.  If it
  is smaller, rearrange the mean deficit and use the exact dyadic harmonic
  lower bound to obtain endpoint growth.
\end{proof}

\begin{theorem}[Terminating maximal-dyadic scan]
  \label{thm:global-scan}
  \lean{GDLowerBound.FourBlock.dyadicScan_normalizedFunctional}
  \leanok
  Starting from any cutoff rank with a functional contribution, the scan
  bounds $\mathcal F(h)$ below by a positive constant divided by
  $B(h)(r(h)+1)^{\beta_0}$.
\end{theorem}

\begin{proof}
  \leanok
  \uses{thm:boundary-propagation,thm:dyadic-dichotomy}
  Choose a maximal dyadic window.  The growth branch is propagated to the
  terminal rank within a fixed factor.  The functional branch moves to a
  strictly later rank, so recursion on the remaining rank gap terminates and
  does not accumulate a loss at every successful branch.
\end{proof}

\begin{theorem}[Exact $1.8288$ normalized floor]
  \label{thm:sharper-floor}
  \lean{GDLowerBound.FourBlock.sharperNormalizedLowerBound}
  \leanok
  There exists $c>0$ such that every nonnegative schedule satisfies the
  normalized schedule floor at the exact exponent $p_0=1.8288$.
\end{theorem}

\begin{proof}
  \leanok
  \uses{thm:bounded-rank,thm:exact-envelope-small-state,thm:global-scan}
  Start the scan at the last small reciprocal-mass state when one exists.  If
  none exists, start from the bounded-rank prefix.  The exact-envelope cutoff
  supplies the cutoff conditions needed by the global scan.
\end{proof}

\begin{theorem}[Sharper gradient-descent lower bound]
  \label{thm:sharper-main}
  \lean{GDLowerBound.FourBlock.sharperMainTheorem}
  \leanok
  For every $1.8288<p<2$ there is a constant $c>0$ such that, for every
  horizon, positive smoothness and radius parameters, and predetermined
  nonnegative schedule, a smooth convex gradient-descent instance satisfies
  \[
    f(x_T)-f_*\geq cLR^2(T+1)^{-p}.
  \]
  Its dimension lies between $1$ and $T+1$, and the initial distance to a
  minimizer is exactly $R$.
\end{theorem}

\begin{proof}
  \leanok
  \uses{thm:floor-reduction,cor:functional-attainment,thm:sharper-floor}
  Monotonicity extends the exact floor to every larger exponent.  Convert it
  to a horizon-rate functional floor, realize the functional geometrically,
  and scale the normalized objective to arbitrary $L$ and $R$.
\end{proof}
